\exercisesection{2}

\begin{enumerate}
\def\labelenumi{\arabic{enumi}.}
\item
  给出一个单射环同态 \(\rho: \mathbf{A} \to \mathbf{B}\) 的例子，使 \(\mathbf{B}\) 是有限 A-代数，映射 \(^a\rho: \operatorname{Spec}(\mathbf{B}) \to \operatorname{Spec}(\mathbf{A})\) 是满射，但 \(\mathbf{B}\) 不是平坦 A-模。反之，给出一个忠实平坦 A-代数的例子，它是有限生成的，但不在 \(\mathbf{A}\) 上整。
\item
  设 A 是使 \(\text{Spec}(A)\) 为 Hausdorff 的环（第二章 § 4 习题 16）。证明：对每个在 A 上整的 A-代数 B，\(\text{Spec}(B)\) 都是 Hausdorff 的。
\item
  设 A 是环，A' 是包含 A 且在 A 上整的环。证明：对 A 的每个理想 a，a 的根等于 A 与 aA' 的根之交。并证明后者是由满足下述条件的元素 \(x \in A'\) 组成的集合：其整依赖方程的系数（首项系数除外）都属于 a（注意，若 y、z 是 A' 的满足这样一个整依赖方程的两个元素，则 y - z 也满足）。
\item
  设 A 是 Noether 环，\(\rho: A \to B\) 是使 B 成为有限生成 A-模的环同态，N 是有限生成 B-模。证明：诸素理想 \(\mathfrak{p} \in \operatorname{Ass}(\rho_{*}(N))\) 恰是在 \(\rho\) 之下诸素理想 \(\mathfrak{q} \in \operatorname{Ass}(N)\) 的原像。
\item
  设 K 是域；在两个未定元的多项式整环 K{[}X, Y{]} 中考虑子环 A = K{[}X \(^{4}\) , Y \(^{4}\) {]}，A' = K{[}X \(^{4}\) , X \(^{3}\) Y, XY \(^{3}\) , Y \(^{4}\) {]}；A 是 Noether 的，A' 是有限 A-代数。设 p = AY \(^{4}\)，它是 A 的一个素理想；证明 m = A'X \(^{4}\) + A'X \(^{3}\) Y + A'XY \(^{3}\) + A'Y \(^{4}\) 是与 A' 的理想 pA' 相伴的一个（浸入）素理想，但并不位于 p 之上（注意 m 是 X \(^{2}\) Y \(^{6}\) 的剩余类在环 A'/pA' 中的零化子）（参见 § 1 习题 26(d)）。
\item
  定义代数数域的一个递增序列 \(K_{0} = Q, K_{1}, \ldots, K_{n}, \ldots\)，使得：若 \(A_{n}\) 表示 Z 在 \(K_{n}, p\) 中的整闭包，其中 p 是一个素数 \(\neq 2\)，则 \(K_{n}\) 是 \(K_{n-1}\) 的二次扩张，且在 \(A_{n}\) 中存在位于 (p) 之上的一个素理想 \(p_{n}\)，使得在 \(A_{n+1}\) 中有两个不同素理想 \(p_{n+1}, p'_{n+1}\) 位于 \(p_{n}\) 之上（注意，在特征 \(\neq 2\) 的有限域中总存在不是平方的元素）。设 K 是诸 \(K_{n}\) 的并，且 \(A'\) 是 Z 在 K 中的整闭包；证明 \(A'\) 中存在无穷多个位于 Z 的极大理想 (p) 之上的（极大）理想。
\item
  设 A 是 Noether 整环，A' 是它的整闭包。证明：对 A 的每个素理想 p，位于 p 之上的 A' 的素理想只有有限多个。（先化归为 A 是以 p 为极大理想的局部环的情形；考虑它的完备化 \(\hat{A}\) 与 \(\hat{A}\) 的全分式环 B，并注意 A 的分式域 K 可等同于 B 的一个子环（第三章 § 3 第 4 节定理 3 的推论 2）。注意 B 的谱是有限的，因而 B 中只有有限多个两两正交的幂等元。然后设 C \(\supset A\) 是 K 的一个子环，它是有限生成 A-模，从而是半局部的。注意 C 的完备化 \(\hat{C}\) 可等同于 B 的一个子环；最后注意到 \(\hat{C}\) 是有限多个局部环的乘积，其个数等于 C 的极大理想的个数，即可得出结论。）把上述结果推广到 Noether 环在其全分式环中的整闭包（化归为约化环的情形）。
\item
  若局部环 A 的整闭包是局部环，则称 A 为单分支的。
\end{enumerate}

\begin{enumerate}
\def\labelenumi{(\alph{enumi})}
\item
  设 A 是局部整环，K 是它的分式域。为使 A 是单分支的，必须且只需 K 的每个作为有限生成 A-模的子 A-代数都是局部环。
\item
  证明：若 Noether 局部环 A 的完备化是整环，则 A 是单分支的（若 B 是包含于 K 中的有限 A-代数，证明由 B 与 \(\hat{A}\) 生成的 \(\hat{A}\) 的分式域的子环同构于 \(\mathbf{B} \otimes_{\mathbf{A}} \hat{\mathbf{A}}\)，为此利用 A-模 \(\hat{\mathbf{A}}\) 的平坦性以及 \(\mathbf{B}\) 包含于某个含于 \(\mathbf{K}\) 的有限生成自由 A-模之中这一事实；然后利用 (a) 及第三章 § 2 第 13 节命题 19 的推论）。
\end{enumerate}

\begin{enumerate}
\def\labelenumi{\arabic{enumi}.}
\setcounter{enumi}{8}
\tightlist
\item
  设 \(k_0\) 是域，\(A'\) 是无限未定元序列的多项式环 \(k_0[X_n]_{n \in \mathbb{N}}\)；它是一个整闭整环，其分式域为 \(K' = k_0(X_n)_{n \in \mathbb{N}}\)。设 \(\sigma\) 是一个 \(k_0\) -自同构，定义在 \(A'\) 上，满足
\end{enumerate}

\[
\sigma (\mathrm{X} _ {2 n}) = - \mathrm{X} _ {2 n} + \mathrm{X} _ {2 n + 1}, \quad \sigma (\mathrm{X} _ {2 n + 1}) = \mathrm{X} _ {2 n + 1}
\]

其中 \(n \geqslant 0\) 任意；\(\sigma\) 与恒等自同构构成 A' 的一个自同构群 G。设 A 是 G 的不变量环，K 是它的分式域；K' 是 K 的次数为 2 的可分扩张，而 A' 是 A 在 K' 中的整闭包。

\begin{enumerate}
\def\labelenumi{(\alph{enumi})}
\tightlist
\item
  证明 \(\mathbf{A}'\) 是一个具有无限基的 A-模。（若 \(k_0\) 的特征 \(\neq 2\)，证明当记 \(\mathbf{Y}_n = \mathbf{X}_{2n} - \frac{1}{2}\mathbf{X}_{2n+1}\) 时，\(\mathbf{A}\) 是 \(\mathbf{A}'\) 中由诸 \(\mathbf{X}_{2n+1}\) 与诸乘积 \(\mathbf{Y}_n\mathbf{Y}_m\)（\(n \geqslant 0\)，\(m \geqslant 0\)）生成的子环。若 \(k_0\) 的特征为 2，则 \(\mathbf{A}'\) 由诸 \(\mathbf{X}_{2n+1}\) 与
\end{enumerate}

\[
\mathrm{X} _ {2 n} ^ {2} + \mathrm{X} _ {2 n} \mathrm{X} _ {2 n + 1}
\]

（\(n\geqslant 0.\)）生成。

\begin{enumerate}
\def\labelenumi{(\alph{enumi})}
\setcounter{enumi}{1}
\item
  设 \(\mathfrak{m}'\) 是 \(\mathbf{A}'\) 的由所有 \(\mathbf{X}_n\) 生成的极大理想，并设 \(\mathfrak{m} = \mathbf{A} \cap \mathfrak{m}'\)；证明 \(\mathfrak{m}' = \mathfrak{mA}'\)，\(\mathfrak{m}'\) 是 \(\mathbf{A}'\) 中位于 \(\mathfrak{m}\) 之上的唯一理想，且域 \(\mathbf{A} / \mathfrak{m}\) 与 \(\mathbf{A}' / \mathfrak{m}'\) 典范同构。
\item
  设 \(\mathfrak{p}'\) 是 \(\mathbf{A}'\) 的由诸 \(\mathbf{X}_{2n+1}\)（\(n \geqslant 0\)）生成的素理想，并设 \(\mathfrak{p} = \mathfrak{p}' \cap \mathbf{A}\)；设 \(k, k'\) 分别是 \(\mathbf{A}/\mathfrak{p}\) 与 \(\mathbf{A}'/\mathfrak{p}'\) 的分式域；\(\mathfrak{p}'\) 是 \(\mathbf{A}'\) 中位于 \(\mathfrak{p}\) 之上的唯一素理想。若 \(k_0\) 的特征不是 2，则 \(k'\) 是 \(k\) 的次数为 2 的可分扩张；若 \(k_0\) 的特征为 2，则 \(k'\) 是 \(k\) 的无限根式扩张。
\end{enumerate}

\begin{enumerate}
\def\labelenumi{\arabic{enumi}.}
\setcounter{enumi}{9}
\tightlist
\item
  \begin{enumerate}
  \def\labelenumii{(\alph{enumii})}
  \tightlist
  \item
    设 \(\mathbf{A}'\) 是两个未定元的多项式环 \(\mathbf{Z}[\mathbf{X},\mathbf{Y}]\)，\(\sigma\) 是 \(\mathbf{A}'\) 的保持 \(\mathbf{Y}\) 不变且满足 \(\sigma (\mathbf{X}) = -\mathbf{X}\) 的自同构；\(\sigma\) 与恒等自同构构成一个自同构群 \(\mathcal{G}\)，作用于 \(\mathbf{A}'\)，\(\mathcal{G}\) 的不变量环为 \(\mathbf{A} = \mathbf{Z}[\mathbf{X}^2,\mathbf{Y}]\)。设 \(\mathfrak{p}'\) 为素理想 \(\mathrm{A}'(\mathrm{Y} - 2\mathrm{X})\)，\(\mathfrak{p} = \mathrm{A}\cap \mathfrak{p}'\)；分解群 \(\mathcal{G}^{\mathbb{Z}}(\mathfrak{p}')\) 即单位元。证明 \(\mathrm{A / p}\) 与 \(\mathrm{A}' / \mathfrak{p}'\) 不同构（考虑这两个环与 \(\mathbf{Z} / 2\mathbf{Z}\) 的张量积）。
  \end{enumerate}
\end{enumerate}

\begin{enumerate}
\def\labelenumi{(\alph{enumi})}
\setcounter{enumi}{1}
\tightlist
\item
  设 K 是特征 \(\neq2\) 的域，A' 是多项式环 K{[}X, Y{]}，G 是 A' 的自同构群，由恒等映射与 A' 的 K{[}Y{]}-自同构 \(\sigma\)（满足 \(\sigma(\mathbf{X}) = -\mathbf{X}\)）组成；G 的不变量环 A 为 \(K[X^{2}, Y]\)。设 \(p'\) 是 A' 的素理想 \(\mathbf{A}'(\mathbf{X}^{3} - \mathbf{Y})\)，\(p = A \cap p'\)；分解群 \(\mathcal{G}^{\mathrm{Z}}(p')\) 即单位元，但 A/p 与 \(A'/p'\) 不同构。
\end{enumerate}

¶ 11. (a) 设 A 是局部整闭整环，m 是它的极大理想，K 是它的分式域，K' 是 K 的有限 Galois 扩张，A' 是 A 在 K' 中的整闭包，m' 是 A' 的位于 m 之上的一个极大理想，B = \(\mathbf{A}^{\mathbb{Z}}(\mathfrak{m}')\) 是它的分解环，n = m' ∩ B。对每个整数 k \textgreater{} 0，有 B = A + n\^{}k（按命题 4 的方式进行论证）。此外，存在 \(\mathbf{K}^{\mathbf{Z}}\)（B 的分式域）的一个本原元 z 使得 z ∈ n；由此推出 n ∩ A{[}z{]} = mA{[}z{]}。

\begin{enumerate}
\def\labelenumi{(\alph{enumi})}
\setcounter{enumi}{1}
\item
  证明：对每个整数 \(k\)，有 \(\mathfrak{m}^k\mathrm{B}_{\mathfrak{n}} \cap \mathrm{A} = \mathfrak{m}^k\)。（若 \(x \in \mathfrak{m}^k\mathrm{B}_{\mathfrak{n}} \cap \mathrm{A}\)，注意存在 \(t \in \mathrm{A}[z]\)（其中 \(z\) 如 (a) 中所定义），使得 \(t \notin \mathfrak{n}\) 且 \(tx \in \mathfrak{m}^k\mathrm{A}[z]\)；记 \(t = t_0 + t_1z + \dots + t_{r-1}z^{r-1}\)（其中 \(r\) 是次数 \([\mathbf{K}^{\mathbf{Z}}: \mathbf{K}]\)），由此推出 \(t_0x \in \mathfrak{m}^k\) 且 \(t_0 \notin \mathfrak{m}\)。）
\item
  现在设 \(\mathbf{K}'\) 是 \(\mathbf{K}\) 的任意 Galois 扩张。保持与上面相同的记号，证明仍有 \(\mathbf{B} = \mathbf{A} + \mathfrak{n}^k\) 与 \(\mathfrak{m}^k\mathrm{B}_{\mathfrak{n}} \cap \mathbf{A} = \mathfrak{m}^k\)（注意 \(\mathbf{A}^{\mathbb{Z}}(\mathfrak{m}')\) 的每个元素都包含在 \(\mathfrak{m}' \cap \mathbf{C}\) 的分解环中，其中 \(\mathbf{C}\) 是 \(\mathbf{A}\) 在 \(\mathbf{K}'\) 的一个有限 Galois 子扩张中的整闭包）。
\item
  在 (c) 的假设下，证明：若 A 是 Noether 的，则 \(B_{n}\) 亦然。（注意 \(B_{n}\) 关于 m-进拓扑的 Hausdorff 完备化依 (c) 可等同于 A 的完备化 \(\hat{A}\)；若 \(\phi: B_{n} \rightarrow \hat{A}\) 是
\end{enumerate}

典范同态，证明：对每个理想 \(a\)（\(B_n\) 中的任意理想），有 \(\phi(a\hat{A}) = a\)，为此利用 \(\hat{A}\) 中每个理想都是有限生成的这一事实，从而 \(a\hat{A}\) 由 \(a\) 的有限多个元素生成，于是化归为 \(K'\) 是 \(K\) 的有限扩张的情形。）

\begin{enumerate}
\def\labelenumi{\arabic{enumi}.}
\setcounter{enumi}{11}
\tightlist
\item
  设 K 是域，\(K'\) 是 K 的有限 Galois 扩张，C 是多项式环 \(K'[X, Y]\)，B 是商环 C/CXY，x、y 是 X、Y 在 B 中的典范像。设 \(A'\) 是 B 的子环 \(K[x] + yK'[y]\)。\(K'\) 在 K 上的 Galois 群 G 忠实作用于 \(A'\)，G 在 \(A'\) 中的不变量环 A 为 \(K[x, y]\)。设 S 是 \(A'\) 中 x 的幂组成的集。证明 \(S^{-1}A' = S^{-1}A\)，且 G 并不忠实作用于 \(S^{-1}A'\)。
\end{enumerate}

¶ 13. (a) 设 K 是特征 0 的域，n 是多项式环 K{[}X, Y, Z{]} 的由 \(Y^{2} - X^{2} - X^{3}\) 生成的理想。证明 n 是素的，且整环 A = K{[}X, Y, Z{]}/n 不是整闭的：若 x、y、z 是 X、Y、Z 在 A 中的像，则 A 的分式域 L 中的元素 t = y/x 在 A 上整，但不属于 A。设 \(A' \supset A\) 是 L 的子环 K{[}x, t, z{]}，它在 A 上整。在 A 中考虑由 xz - y 与 \(z^{2} - 1 - x\) 生成的素理想 p，以及由 x、y、z - 1 生成的极大理想 \(q \supset p\)。另一方面，在 A' 中考虑极大理想 \(q'\)，它由 x、\(t + 1\) 与 z - 1 生成；证明 \(q'\) 位于 q 之上，但不存在 A' 的位于 p 之上的素理想 \(p' \subset q'\)。（考虑一个从 \(A'/pA'\) 到 K 的扩域的同态，并证明在这样的同态下 \(q'\) 的像必为 0。）

\begin{enumerate}
\def\labelenumi{(\alph{enumi})}
\setcounter{enumi}{1}
\item
  设 \(\mathbf{A}'\) 是多项式环 \(\mathbf{Z}[\mathbf{X}]\) 对理想 \(n\)（由 2X 与 \(\mathbf{X}^2 -\mathbf{X}\) 生成）的商环；设 \(x\) 是 \(\mathbf{X}\) 在 \(\mathbf{A}'\) 中的像；\(\mathbf{A}'\) 在 \(\mathbf{Z}\) 上整，但数 2 是 \(\mathbf{A}'\) 中的零因子。设 \(q'\) 是 \(\mathbf{A}'\) 的由 2 与 \(x - 1\) 生成的素理想；则 \(q' \cap \mathbf{Z} = 2\mathbf{Z}\)；若 \(p\) 是 \(\mathbf{Z}\) 中的素理想 (0)，证明不存在素理想 \(p' \subset q'\)（含于 \(\mathbf{A}'\) 且位于 \(p\) 之上）。
\item
  设 K 是域，n 是多项式环 K{[}X, Y, Z{]} 的由 \(X^{2}\)、XY、XZ、YZ - Y 与 \(Z^{2} - Z\) 生成的理想，A' 是商环 K{[}X, Y, Z{]}/n，x、y、z 是 X、Y、Z 在 A' 中的像，A 是 A' 的子环 K{[}x, y{]}；A 在它的全分式环中整闭，而 A' 在 A 上整。在 A 中考虑由 x 生成的素理想 p 以及由 x 与 y 生成的极大理想 q。另一方面，在 A' 中考虑由 x、y、z 生成的极大理想 \(q'\)；证明 \(q'\) 位于 q 之上，但不存在位于 p 之上的素理想 \(p' \subset q'\)（方法同 (a)）。
\end{enumerate}

* 试从几何上解释 (a) 与 (c) 中的例子。 *

¶ 14. (a) 设 A 是域 K 的子环，x 是 K 中一个 \(\neq0\) 的元素。证明：若 x 不在 A 上整，则存在 \(A[x^{-1}]\) 的一个包含 \(x^{-1}\) 的极大理想，且对 \(A[x^{-1}]\) 的每个包含 x 的极大理想 m，理想 \(A\cap m\) 都是 A 的极大理想（注意 x 不属于环 \(A[x^{-1}]\)）。

\begin{enumerate}
\def\labelenumi{(\alph{enumi})}
\setcounter{enumi}{1}
\item
  设 A 是局部整环，m 是它的极大理想，K 是包含 A 的域，k 是 A 的剩余域。设 \(x \in K\) 满足 \(x \notin A\) 且 \(x^{-1} \notin A\)，并设 A 在环 A{[}x{]} 中整闭。证明存在一个 A-代数同态 \(A[x] \to k[X]\)，它延拓典范同态 \(A \to k\)，且把 x 映为 X。（需要证明：若 \(\sum_{i=0}^{n}a_{i}x^{i}=0\)，其中对所有 i 有 \(a_{i}\in A\)，则所有 \(a_{i}\) 都属于 m。否则将有一个关系式 \(a_{n}x^{n-p}+\cdots+a_{p}=-(a_{p+1}x^{-1}+\cdots+a_{0}x^{-p})\)，其中 \(a_{p}\) 在 A 中可逆；利用 §1 习题 5，证明两端的公共值 b 属于 A，再利用 (a)，证明必有 \(b\in m\)；由此推出 \(x^{-1}\) 将在 A 上整，这是荒谬的。）
\item
  在 (b) 的假设下，证明 A{[}x{]} 的理想 mA{[}x{]} 是素的且非极大的。
\end{enumerate}

\begin{enumerate}
\def\labelenumi{\arabic{enumi}.}
\setcounter{enumi}{14}
\item
  设 A 是环，p 是 A 的素理想。证明：在多项式环 B = A{[}X{]} 中，理想 pB = q 是素的；局部环 \(B_{q}\) 的极大理想等于 \(pB_{q}\)；若 k 是 A/p 的分式域，则 \(B_{q}/pB_{q}\) 同构于有理分式域 \(k(\mathbf{X})\)。若 A 是整闭的，则 \(B_{q}\) 亦然。
\item
  设 A 是整环，K 是它的分式域，\(K'\) 是 K 的次数为 n 的有限扩张，\(A'\) 是 \(K'\) 中包含 A 的子环，以 \(K'\) 为分式域且在 A 上整。设 m 是 A 的一个极大理想，B 是局部环 \(A[X]_{mA[X]}\)（习题 15）。
\end{enumerate}

\begin{enumerate}
\def\labelenumi{(\alph{enumi})}
\item
  子环 \(\mathbf{B}' = \mathbf{B}[\mathbf{A}']\) 含于 \(\mathbf{K}'(\mathbf{X})\)，它在 \(\mathbf{B}\) 上整；它的分式域是 \(\mathbf{K}'(\mathbf{X})\)，且 \([\mathbf{K}'(\mathbf{X}) : \mathbf{K}(\mathbf{X})] = n\)。
\item
  设 \(\mathfrak{m}'\) 是 \(\mathbf{A}'\) 的包含 \(\mathfrak{m}\) 的任一极大理想，则 \(\mathfrak{n}' = \mathfrak{m}'\mathbf{B}'\) 是 \(\mathbf{B}'\) 的极大理想；若记 \(\mathfrak{n} = \mathfrak{m}\mathbf{B}\)，则
\end{enumerate}

\[
[ (\mathbf {A} ^ {\prime} / \mathfrak {m} ^ {\prime}): (\mathbf {A} / \mathfrak {m}) ] = [ (\mathbf {B} ^ {\prime} / \mathfrak {n} ^ {\prime}): (\mathbf {B} / \mathfrak {n}) ]
\]

以及

\[
[ (\mathrm{A} ^ {\prime} / \mathfrak {m} ^ {\prime}): (\mathrm{A} / \mathfrak {m}) ] _ {s} = [ (\mathrm{B} ^ {\prime} / \mathfrak {n} ^ {\prime}): (\mathrm{B} / \mathfrak {n}) ] _ {s}.
\]

（注意 \((\mathbf{A}' / \mathfrak{m}') \otimes_{\mathbf{A}} \mathbf{B}\) 同构于域 \((\mathbf{A}' / \mathfrak{m}')(\mathbf{X})\)，由此推出 \(\mathbf{B}' / \mathfrak{n}'\) 同构于 \((\mathbf{A}' / \mathfrak{m}')(\mathbf{X}).\)）

\begin{enumerate}
\def\labelenumi{(\alph{enumi})}
\setcounter{enumi}{2}
\tightlist
\item
  映射 \(m' \mapsto m'B'\) 是从 \(A'\) 的包含 m 的极大理想所成之集到 \(B'\) 的极大理想所成之集上的双射；逆双射是 \(n' \mapsto n' \cap A'\)。
\end{enumerate}

¶ 17. (a) 设 K 是无限域，E 是 K 的有限可分代数扩张，g 是 K{[}X{]} 的一个多项式。证明存在元素 \(x \in E\)，使得 \(\mathrm{E} = \mathrm{K}(x)\)，且 x 在 K 上的极小多项式与 g 互素（按《代数》第五章 § 7 第 7 节命题 12 的方式论证）。

\begin{enumerate}
\def\labelenumi{(\alph{enumi})}
\setcounter{enumi}{1}
\item
  设 \(k\) 是域，A 是准素 \(k\) -代数（第四章 § 2 习题 32），\(m\) 是它的唯一素理想。若 \(A / m\) 是 \(k\) 的超越扩张，则 A 中存在不在 \(k\) 上代数的元素。若 \(A / m\) 在 \(k\) 上代数，且 \(d\) 是一个至少等于 \(A / m\) 在 \(k\) 上的次数的可分部分的整数（从而当该可分部分为无限时，d 可为任意整数），证明 A 中存在在 \(k\) 上的极小多项式的次数 \(\geqslant d\) 的元素。进而，若 \(A / m\) 在 \(k\) 上不可分，或 \(m \neq 0\)，则 A 中存在一个元素，它在 \(k\) 上的极小多项式的次数 \(>d\)。（若 \(A / m\) 在 \(k\) 上不可分，利用《代数》第五章 § 8 习题 5。若 \(A / m\) 的次数为 \(d\) 且在 \(k\) 上可分，又 \(x \in A\) 使 \(\xi \in A / m\)（即 \(x\) 的剩余类）是 \(A / m\) 的本原元，其极小多项式为 \(f \in k[X]\)，则注意 \(f'(\xi) \neq 0\)，并由此推出：若 \(y \neq 0\) 是 \(m\) 的一个元素，则 \(f(x + y) \neq 0\)。）
\item
  设 \(k\) 是无限域，\(A\) 是 \(k\) -代数，为 \(r\) 个准素代数 \(A_{i} (1 \leqslant i \leqslant r)\) 的直合成；设 \(m_{i}\) 是 \(A_{i}\) 的极大理想，\(d_{i}\) 是一个整数，它至少等于 \(A_{i} / m_{i}\) 在 \(k\) 上的次数的可分部分（当 \(A_{i}\) 代数且该部分有限时），否则为任意整数。证明 \(A\) 中存在一个不在 \(k\) 上代数的元素，或者
\end{enumerate}

其在 \(k\) 上的极小多项式的次数 \(\geqslant \sum_{i=1}^{r} d_i\)；进而，若某个 \(m_i\) 为 \(\neq 0\)，或某个域 \(A_i / m_i\) 不是 \(k\) 的可分代数扩张，则 \(A\) 中存在一个不在 \(k\) 上代数的元素，或者其极小

多项式在 \(k\) 上的次数 \(> \sum_{i=1}^{r} d_i\)（利用 (a) 与 (b)）。

¶ 18. 设 A 是局部整闭整环，K 是它的分式域，m 是它的极大理想，\(k = \mathbf{A} / \mathfrak{m}\) 是它的剩余域。设 \(\mathbf{K}'\) 是 \(\mathbf{K}\) 的次数为 \(n\) 的有限代数扩张，\(\mathbf{A}'\) 是 \(\mathbf{K}'\) 中包含 \(\mathbf{A}\) 的子环，它在 \(\mathbf{A}\) 上整且以 \(\mathbf{K}'\) 为分式域，\(\mathfrak{m}_i\)（\(1 \leqslant i \leqslant r\)）是 \(\mathbf{A}'\) 的诸极大理想。称理想 \(\mathfrak{m}\) 在 \(\mathbf{A}'\) 中是非分歧的，如果存在 \((w_i)_{1 \leqslant i \leqslant n}\)，它构成 \(\mathbf{K}'\) 在 \(\mathbf{K}\) 上的一组基，由 \(\mathbf{A}'\) 的元素组成，且其判别式 \(D_{\mathbf{K}' / \mathbf{K}}(w_1, \ldots, w_n)\)（《代数》第九章 §2）属于 \(\mathbf{A} - \mathfrak{m}\)。

\begin{enumerate}
\def\labelenumi{(\alph{enumi})}
\item
  证明：若 \(\mathfrak{m}\) 在 \(\mathbf{A}'\) 中非分歧，则 \(\mathbf{K}'\) 是 \(\mathbf{K}\) 的可分扩张，\(\mathbf{A}'\) 是 \(\mathbf{A}\) 在 \(\mathbf{K}'\) 中的整闭包，并且是自由 \(\mathbf{A}\) -模，它在 \(\mathbf{A}\) 上的任何基都是 \(\mathbf{K}'\) 在 \(\mathbf{K}\) 上的基；\(\mathfrak{mA}'\) 是 \(\mathbf{A}'\) 的 Jacobson 根，而 \(\mathbf{A}' / \mathfrak{mA}'\) 是 \(k\) 上秩为 \(n\) 的半单 \(k\) -代数，它在 \(k\) 上可分，且是诸域 \(\mathbf{A}' / \mathfrak{m}_i'\) 的直合成。
\item
  设 \(d_{i}\) 是 \(A'/m_{i}'\) 在 \(A/m\) 上的次数的可分部分。证明：若 \(k\) 是无限的，则 \(\sum_{i=1}^{r} d_{i} \leqslant n\)（利用习题 17 (c)）；此外，下列条件等价：
\end{enumerate}

（α）m 在 A' 中非分歧。

\((\beta)\) \(\mathbf{A}'\) 是有限生成 A-模，且 \(\mathbf{A}' / \mathfrak{mA}'\) 是半单 \(k\) -代数并在 \(k\) 上可分。

\((\gamma)\sum_{i=1}^{r}d_{i}\leqslant n.\)

（为看出 \((\beta)\) 蕴涵 \((\gamma)\)，注意，利用第二章 § 3 第 2 节命题 4 的推论 2，A-模 \(A'\) 有一个生成系，其元素模 \(mA'\) 的剩余类构成 \(A'/mA'\) 在 \(k\) 上的一组基；由此推出

\[
[ \mathbf {K} ^ {\prime}: \mathbf {K} ] \leqslant [ (\mathbf {A} ^ {\prime} / \mathsf {m A} ^ {\prime}): k ].
\]

为看出 \((\gamma)\) 蕴涵 \((\alpha)\)，利用习题 17 (c) 证明：存在 \(\mathbf{A}' / \mathfrak{mA}'\) 在 \(k\) 上的一组基，由一个元素的诸幂 \(\xi^i\)（\(\xi\) 为该元素，\(0 \leqslant i \leqslant n - 1\)）组成；若 \(x \in \mathbf{A}'\) 是剩余类 \(\xi\) 的一个元素，则推出 \(\mathrm{D}_{\mathrm{K}' / \mathrm{K}}(1, x, \ldots, x^{n-1})\) 属于 \(\mathbf{A} - \mathfrak{m}\)。）给出一个例子，其中 \(\mathbf{K}'\) 在 \(\mathbf{K}, \mathbf{A}' / \mathfrak{mA}' = \mathbf{A} / \mathfrak{m}\) 上可分，而 \(\mathbf{A}'\) 不是有限生成 A-模（参见习题 9 (c)）。

\begin{enumerate}
\def\labelenumi{(\alph{enumi})}
\setcounter{enumi}{2}
\tightlist
\item
  把 (b) 的结果推广到 \(k\) 为有限的情形（在习题 16 的记号下，证明：若 \(n\) 在 \(B'\) 中非分歧，则 \(m\) 在 \(A'\) 中非分歧；利用 \(B / n\) 是无限的这一事实，使用习题 16 (c)，从而得到一个元素 \(z \in B'\)，使得
\end{enumerate}

\[
\mathrm{D} _ {\mathbf {K} ^ {\prime} (\mathbf {X}) / \mathbf {K} (\mathbf {X})} (1, z, \dots , z ^ {n - 1}) \in \mathrm{B} - n;
\]

证明可设 \(z\) 是 \(\mathbf{A}'[\mathbf{X}]\) 中的一个多项式，并且把上面的判别式表为 \(\mathbf{X}\) 的多项式，便得到由 \(n\) 个元素 \(w_{i}\)（\(1 \leqslant i \leqslant n\)）组成的一个组，属于 \(\mathbf{A}'\)，使得 \(\mathrm{D}_{\mathrm{K}'/\mathrm{K}}(w_1, \ldots, w_n) \in \mathbf{A} - \mathfrak{m}\)）。

¶ 19. 设 A 是整闭整环，K 是它的分式域，K' 是 K 的次数为 n 的有限代数扩张，A' 是 K' 中包含 A 的子环，它在 A 上整且以 K' 为分式域；若 p 是 A 的素理想，用 \(\mathbf{A}_{\mathfrak{p}}^{\prime}\) 表示 \(\mathbf{A}^{\prime}\) 的分母属于 \(\mathbf{A}-\mathfrak{p},\mathfrak{p}\) 的分式环；则称 p 在 \(\mathbf{A}^{\prime}\) 中非分歧，如果 \(\mathfrak{p}\mathbf{A}_{\mathfrak{p}}\) 在 \(\mathbf{A}_{\mathfrak{p}}^{\prime}\) 中非分歧。

\begin{enumerate}
\def\labelenumi{(\alph{enumi})}
\tightlist
\item
  设 \(\mathfrak{p}_i^{\prime}(1\leqslant i\leqslant r)\) 是 \(\mathbf{A}'\) 的位于 \(\mathfrak{p}\) 之上的诸素理想。设 \(d_{i}\) 是 \(\mathbf{A}' / \mathfrak{p}_i'\) 的分式域在
\end{enumerate}

A/p 的分式域上的次数的可分部分；证明 \(\sum_{i=1}^{r} d_i \leqslant n\)。（化归为习题 18。）

\begin{enumerate}
\def\labelenumi{(\alph{enumi})}
\setcounter{enumi}{1}
\tightlist
\item
  下列条件等价：
\end{enumerate}

（α）理想 p 在 A' 中非分歧。

（β）A' 包含 K' 在 K 上的一组基，其判别式属于 A - p。

\((\gamma)\sum_{i=1}^{r}d_{i}=n.\)

特别地，若 \(r = n\)，则 \(\mathfrak{p}\) 在 \(\mathbf{A}'\) 中非分歧。

给出一个例子，其中 \(p\) 在 \(A'\) 中非分歧，但 \(A'\) 不是有限生成 \(A\) -模（参见习题 9 (a)）。

\begin{enumerate}
\def\labelenumi{(\alph{enumi})}
\setcounter{enumi}{2}
\item
  若 A 的每个素理想都在 A' 中非分歧，则称 A 在 A'（或在 K'）中非分歧。证明此时 A' 是忠实平坦 A-模（参见第二章 § 3 第 4 节命题 15 的推论）。设 K'' 是 K' 的次数为 n' 的有限代数扩张；设 A' 是 A 在 K' 中的整闭包，A'' 是 A 在 K'' 中的整闭包；为使 A 在 A'' 中非分歧，必须且只需 A 在 A' 中非分歧且 A' 在 A'' 中非分歧。
\item
  设 \(k\) 是域，\(A\) 是整闭的整 \(k\) -代数，\(K\) 是它的分式域。设 \(k'\) 是 \(k\) 的有限可分代数扩张，次数为 \(n\)（在 \(k\) 上）。证明：若 \(k' \otimes_k K\) 是域，则 \(A\) 在 \(k' \otimes_k K\) 中非分歧（考虑本原元 \(x\)（\(k'\) 在 \(k\) 上的）以及 \(k' \otimes_k K\) 的由诸 \(x^i\)（\(0 \leqslant i \leqslant n - 1\)）组成的基）。
\end{enumerate}

\begin{enumerate}
\def\labelenumi{\arabic{enumi}.}
\setcounter{enumi}{19}
\item
  设 A 是第三章 § 3 习题 15 (b) 中所定义的整环，C 是局部环 \(\mathbf{A}_{\mathfrak{m}}\)，\(\mathfrak{n} = \mathfrak{mA}_{\mathfrak{m}}\) 是 C 的极大理想，L 是 A 与 C 的分式域。若 \(x, y\) 是 X 与 Y 在 C 中的典范像，证明 \(t = y / x\) 不属于 C，但 \(t^2 \in \mathbf{C}\)，从而 C 不是整闭的；此外 \(\mathbf{C}' = \mathbf{C}[t]\) 是 C 的整闭包，且同构于 \(\mathbf{S}^{-1}\mathbf{K}[\mathbf{T}]\)（T 为未定元），其中 S 是 K{[}T{]} 中由既不被 T - 1 整除、也不被 T + 1 整除的多项式组成的乘性子集。证明 \(\mathbf{C}'\) 是有限生成 C-模，且 \(\mathbf{C}'\) 的位于 \(\mathfrak{n}\) 之上的（极大）理想是主理想 \(\mathfrak{q}_1 = (t - 1)\mathbf{C}'\) 与 \(\mathfrak{q}_2 = (t + 1)\mathbf{C}'\)；但局部环 \(\mathbf{C}_{\mathfrak{q}_1}'\) 与 \(\mathbf{C}_{\mathfrak{q}_2}'\) 不是在 C 上整的 C-代数。
\item
  设 A 是包含域 Q 的整闭整环，K 是它的分式域，L 是 K 的代数扩张，B 是 A 在 L 中的整闭包。证明：对 A 的每个理想 a，有 \(A \cap aB = a\)。
\item
  设 A 是 Noether 整环，K 是它的分式域，A' 是 K 中包含 A 的子环。设 A 在 A' 中整闭，且对 A 的每个素理想 p，存在素理想 \(p'\)（属于 \(A'\)），位于 p 之上。用反证法论证：设 \(x \in A' \cap C A\)，令 \(a = A \cap (x^{-1} A)\)，它是异于 A 的理想。设 p 是与 A/a 相伴的素理想；则传递子 b = a: p 异于 a（第四章 §1 第 4 节命题 9），因而存在 \(y \in b\)，使得 \(xy \notin A\)，且
\end{enumerate}

\[
x y p \subset p A ^ {\prime} \cap A = p;
\]

由此得出矛盾。）

